% =====================================================================
%  1bat-ud2-13.tex  —  HORA 13: Veure les arrels amb GeoGebra
%  (sense preàmbul: s'inclou des de main.tex)
% =====================================================================
\horatitol{Sessió 13 — Veure les arrels amb GeoGebra}%
{Comprovar amb una eina gràfica que les arrels calculades a mà coincideixen amb el tall a l'eix horitzontal.}

\subsection{Idea clau}

A les sessions 10, 11 i 12 hem trobat arrels \textbf{algebraicament} (aïllant
la $x$ o aplicant la fórmula general) i les hem interpretat com el llindar
d'exclusió d'un ajut. Avui canviem d'eina: amb \textbf{GeoGebra}
representarem gràficament aquests mateixos polinomis i comprovarem, amb els
ulls, que l'arrel calculada coincideix exactament amb el punt on la gràfica
\textbf{talla l'eix horitzontal}.

\begin{idea}
\textbf{L'arrel, vista a la gràfica.}
Si $P(x)$ és un polinomi i el representem com una corba $y=P(x)$:
\begin{itemize}[leftmargin=1.4em, itemsep=2pt]
  \item Una \textbf{arrel} de $P(x)$ és un valor de $x$ on la corba
        \textbf{talla l'eix horitzontal} (és a dir, $y=0$).
  \item A l'esquerra del tall, si la corba és per sobre de l'eix
        ($y>0$), l'ajut \textbf{encara existeix}.
  \item A la dreta del tall, si la corba passa per sota de l'eix
        ($y<0$), l'ajut ja \textbf{s'ha exhaurit} (recordeu, de la sessió
        12, que aquest valor negatiu no té sentit real: s'interpreta com
        a $0\eur$).
\end{itemize}
\end{idea}

\subsection{Exemples}

\begin{exemple}[introduir un polinomi a GeoGebra]
Obrim la vista gràfica de GeoGebra i, a la barra d'entrada, escrivim
directament l'ajut de les sessions 10 i 12:
\[
B(x) = -20x + 600
\]
GeoGebra dibuixa a l'instant una recta. Activant la icona d'\emph{intersecció
amb els eixos} (o escrivint \texttt{Arrel(B)} a la barra d'entrada),
l'eina marca el punt exacte on la recta talla l'eix horitzontal i en
mostra les coordenades. Hauríeu de veure el punt $(30, 0)$: exactament
l'arrel que vam calcular a mà a la sessió 10.
\end{exemple}

\begin{exemple}[un polinomi de segon grau: dues arrels]
Introduïm ara el cost net del casal de joves de la sessió 11:
\[
M(x) = x^2 - 14x + 40
\]
La gràfica ja no és una recta sinó una \textbf{paràbola}. Amb
\texttt{Arrel(M)}, GeoGebra marca \textbf{dos} punts de tall:
$(4, 0)$ i $(10, 0)$, coincidint amb les arrels $x_1=4$ i $x_2=10$
calculades amb la fórmula general. Es pot comprovar visualment allò que
ja sabíem algebraicament: \textbf{entre} els dos talls, la paràbola passa
per sota de l'eix (cost net negatiu, és a dir, superàvit), i per fora
dels dos talls, per sobre (cal aportar la diferència).
\end{exemple}

\subsection{Exercicis resolts}

\begin{exresolt}[ajustar la finestra de visualització]
En introduir $C(x)=\num{0,05}x^2-10x+500$ (el copagament de la sessió 4)
a GeoGebra amb la finestra per defecte, no es veu cap tall amb l'eix
horitzontal. Per què, i què cal fer?
\solucio
La finestra per defecte de GeoGebra sol mostrar $x$ i $y$ entre,
aproximadament, $-10$ i $10$. Sabem, de la sessió 11, que l'arrel
d'aquest polinomi és $x=100$ (arrel doble, $\Delta=0$), un valor molt
fora d'aquest rang. A més, $C(0)=500$, un valor de $y$ també molt gran.
Cal \textbf{ajustar la finestra de visualització} (per exemple, $x$ entre
$0$ i $200$, $y$ entre $0$ i $600$) perquè el punt de tall hi sigui
visible. Un cop ajustada, la paràbola es veu \textbf{tangent} a l'eix
horitzontal exactament a $x=100$: toca l'eix en un sol punt sense
travessar-lo, la imatge gràfica d'una arrel doble.
\resposta{Cal ampliar la finestra de visualització (els valors de la
recerca de l'arrel, $x=100$, queden fora del rang per defecte); un cop
ajustada, la paràbola és tangent a l'eix a $x=100$.}
\end{exresolt}

\begin{exresolt}[comprovar un cas sense arrels reals]
Introduïu a GeoGebra el polinomi $D(x)=2x^2-4x+10$ de la sessió 11
(per al qual vam calcular $\Delta=-64<0$). Què hauríeu d'esperar veure
a la gràfica, i com ho interpreteu en el context del servei que
modelitza?
\solucio
Com que $\Delta<0$, l'equació $D(x)=0$ no té cap solució real. A la
gràfica, això es tradueix en una paràbola que \textbf{no talla mai}
l'eix horitzontal: per molt que s'ampliï o es desplaci la finestra de
visualització, la corba sempre queda per sobre de l'eix ($a=2>0$).
En GeoGebra, en demanar \texttt{Arrel(D)}, l'eina no retorna cap punt.
En el context del servei, això confirma que el seu cost net mai
s'anul·la: és sempre positiu, per a qualsevol nombre d'usuaris.
\resposta{La paràbola no talla mai l'eix horitzontal (sempre per sobre);
GeoGebra no troba cap arrel real, coherent amb $\Delta=-64<0$: el cost
net del servei no s'anul·la mai.}
\end{exresolt}

\subsection{Exercicis per fer}

\noindent\textit{Tasques de comprovació visual: introduïu cada polinomi a
GeoGebra i contrasteu el que veieu amb el càlcul algebraic ja fet a les
sessions anteriors.}

\begin{exercicis}
\begin{exlist}
  \item Introdueix a GeoGebra el polinomi $B(x)=-20x+600$ (sessions 10 i
        12). Localitza el punt de tall amb l'eix horitzontal. Coincideix
        amb l'arrel $x=30$ calculada a la sessió 10?

  \item Introdueix el polinomi de segon grau $x^2-7x+10$ (treballat a la
        sessió 11). Observa quants punts de tall té la paràbola amb
        l'eix horitzontal. Coincideixen amb les arrels que vau calcular
        aleshores?

  \item Introdueix el copagament $C(x)=\num{0,05}x^2-10x+500$ (sessions 4
        i 11). Ajusta la finestra de visualització fins que es vegi
        clarament el punt on la paràbola toca l'eix. Anota les
        coordenades d'aquest punt: coincideixen amb l'arrel doble
        $x=100$?

  \item El servei de teleassistència de la sessió 11 es modelitza amb
        $T(x)=\num{0,05}x^2-8x+320$. Introdueix-lo a GeoGebra i localitza
        el seu tall (o talls) amb l'eix horitzontal. Compara el resultat
        gràfic amb el càlcul algebraic que vau fer aleshores.

  \item \textbf{(Compte amb la finestra)} Introdueix el polinomi
        $M(x)=-25x+750$ (sessió 12, servei de menjador social). Si en
        ampliar GeoGebra per defecte no veus cap tall amb l'eix, no
        donis per fet que no n'hi ha: comprova primer si el punt
        simplement queda fora de la finestra visible, com va passar amb
        el copagament de l'exercici resolt 13.1.

  \item \textbf{(Interpretació)} Introdueix el polinomi
        $D(x)=2x^2-4x+10$ (sessió 11, sense arrels reals) i confirma a
        GeoGebra que la paràbola no talla mai l'eix horitzontal. Explica,
        en context, què significaria per a un servei que la seva gràfica
        de cost net tallés l'eix horitzontal en algun punt, i
        què significa que no ho faci mai.
\end{exlist}
\end{exercicis}

\noindent\textbf{Atenció a la diversitat:} si és la primera vegada que useu
GeoGebra, teniu a la pissarra digital la pauta amb captures dels passos
bàsics (introduir una expressió a la barra d'entrada, activar la icona
d'intersecció amb els eixos o escriure \texttt{Arrel(...)}, ajustar la
finestra de visualització amb la lupa o arrossegant el pla). Les parelles
estan formades perquè qui tingui més soltesa digital pugui acompanyar qui
en tingui menys; l'objectiu central segueix sent matemàtic (relacionar
arrel i tall a l'eix), no la destresa amb l'eina.

% ---------------------------------------------------------------------
\fullsolucions{Sessió 13}
\begin{solucions}
\begin{sollist}
  \item Sí: el punt de tall és $(30, 0)$, coincidint amb l'arrel $x=30$
        calculada algebraicament a la sessió 10.
  \item Dos punts de tall, $(2, 0)$ i $(5, 0)$, coincidint amb les arrels
        $x_1=5$, $x_2=2$ calculades amb la fórmula general a la sessió 11.
  \item El punt és $(100, 0)$; la paràbola és tangent a l'eix en aquest
        punt (el toca sense travessar-lo), coherent amb l'arrel doble
        ($\Delta=0$) calculada a la sessió 11.
  \item Un únic punt de tall, $(80, 0)$ (arrel doble, $\Delta=0$),
        coincidint amb el càlcul algebraic de la sessió 11.
  \item El tall amb l'eix és a $(30, 0)$; cal ampliar la finestra de
        visualització (per exemple $x$ entre $0$ i $40$) perquè s'hi
        vegi, ja que amb la finestra per defecte queda fora del rang
        visible.
  \item La paràbola no talla mai l'eix horitzontal (sempre per sobre,
        coherent amb $\Delta=-64<0$). Si la gràfica tallés l'eix,
        voldria dir que existeix algun nombre d'usuaris pel qual el
        cost net del servei és exactament zero (un llindar real); que
        no el talli mai significa que el cost net és sempre positiu,
        per a qualsevol nombre d'usuaris.
\end{sollist}
\end{solucions}
